\section{李群——李代数“辞典”}

李从 1874 年起在大量论文中发展起来的“有限连续”群理论，在宏伟的专著《变换群论》（Theorie der Transformationsgruppen）（{[}203{]}，1888-1893）中得到了系统的阐述，该书是与 F. 恩格尔合作写成的；\(^{7}\) 该理论是第一卷以及第三卷最后五章的对象，第二卷则专用于切触变换。

正如书名所表明的，这部著作中讨论的从来只是变换群，即方程 (25.4) 意义下的变换群，其中“变量”\(x_{i}\) 的空间与“参数”\(a_{j}\) 的空间起初起着同等重要的作用。此外，“抽象”群的概念在当时还没有被清楚地表达出来；1883 年，当李（{[}202{]}，vol.~5, Abh. XII, pp.~311-313）注意到采用 (25.10) 的记号时，方程 \(w = H(u, v)\)——它给出群的两个变换之积的参数——定义了一个新的群，他是把它作为参数空间上的一个变换群来考虑的，由此得到他所说的“参数群”（他甚至得到了两个，它们不是别的，正是左平移群与右平移群）。\(^{8}\)

方程 (25.4) 中的变量 \(x_{i}\) 与参数 \(a_{j}\) 原则上被假定为复的（第三卷第 XIX-XXIV 章除外），函数 \(f_{i}\) 为解析的；李与恩格尔当然意识到，这些函数一般并非对 \(x_{i}\) 与 \(a_{j}\) 的一切复值都有定义，因而这种变换的合成会引起严重的困难（{[}203{]}，v. I, pp.~15-17, pp.~33-40 以及散见各处）；虽然后来他们表述时几乎总是显得他们所研究的变换的合成总是可以无限制地进行，这无疑只是为了叙述的方便，每逢必要他们就明确地重新确立“局部”的观点（例如参见 loc. cit., pp.~168 或 189，或 ibid., v. 3, p.~2, 页末注）；换句话说，他们所研究的数学对象接近于我们今天所说的偏运算律。他们也不忘在适当场合考虑整体群，例如经典群的 4 个系列（{[}203{]}，v. 3, p.~682），但似乎没有给自己提出过如下问题：一般地说一个“整体群”可以是什么；对他们来说，只要能对经典群的“参数”（这些群的“变量”不引起任何困难，因为所涉及的是 \(\mathbf{C}^n\) 的线性变换）得到恒等变换邻域内的“局部”参数系统就够了，而不必关心他们所写下的公式的有效域。然而他们给自己提出了一个明显越出局部理论的问题：⁹ 研究“混合”群，也就是只有有限多个连通分支的群，例如正交群（{[}203{]}，v. I, p.~7）。他们把这项研究表述为对一个对合成和取逆封闭的变换集合的研究，该集合是诸集合 \(H_{j}\) 的并，其中每个集合由形如 (25.4) 的函数组 \(f_{i}^{(j)}\) 所描述；每个 \(H_{j}\) 的（本质）参数的数目甚至被先验地假定依赖于 j，但他们实际上证明这个数目对所有 \(H_{j}\) 都相同。他们的主要结果便是一个有限连续群 G 的存在性，使得 \(H_{j} = G \circ h_{j}\)（对某个 \(h_{j} \in H_{j}\) 以及所有 j）；他们还建立了 G 在混合群中是正规的，并注意到后者的不变量的确定可化归为 G 的不变量与一个不连续群的不变量的确定（{[}203{]}，v. I, chap.~18）。

{[}203{]}中发展起来的一般理论最终（作者们并未以十分系统的方式陈述这一点）锻造出了一部“辞典”，把“有限连续”群的性质翻译为其无穷小变换全体的性质。它建立在“李的三条定理”之上，其中每一条都由一个断言及其逆命题组成。

第一条定理（{[}203{]}，v. I, pp.~33 与 72 以及 v. 3, p.~563）首先断言：若在 (25.4) 中参数是本质的，则函数 \(f_{i}\) 满足形如

\[
\frac {\partial f _ {i}}{\partial a _ {j}} = \sum_ {k = 1} ^ {r} \xi_ {k i} (f (x, a)) \psi_ {k j} (a) (1 \leq i \leq n)\tag{25.15}
\]

其中矩阵 \((\xi_{kl})\) 有最大秩，且 \(\det(\psi_{kj}) \neq 0\)；反过来，若函数 \(f_{i}\) 具有这一性质，则公式 (25.4) 定义一个变换群芽。

第二条定理（{[}203{]}，v. I, pp.~149 与 158，以及 v. 3, p.~590）给出 \(\xi_{ki}\) 之间以及 \(\psi_{ij}\) 之间的关系：关于 \(\xi_{ki}\) 的条件可以写成形式

\[
\sum_ {k = 1} ^ {n} \left(\xi_ {i k} \frac {\partial \xi_ {j l}}{\partial x _ {k}} - \xi_ {j k} \frac {\partial \xi_ {i l}}{\partial x _ {k}}\right) = \sum_ {k = 1} ^ {r} c _ {i j} ^ {k} \xi_ {k l} (1 \leq i, j \leq r, 1 \leq l \leq n)\tag{25.16}
\]

其中 \(c_{ij}^{k}\) 是常数 \((1 \leq i, j, k \leq r)\)，对 i, j 反对称。毛雷尔所给形式的关于 \(\psi_{ij}\) 的条件是：

\[
\frac {\partial \psi_ {k l}}{\partial a _ {m}} - \frac {\partial \psi_ {k m}}{\partial a _ {l}} = \frac {1}{2} \sum_ {1 \leq i, j \leq r} c _ {i j} ^ {k} (\psi_ {i l} \psi_ {j m} - \psi_ {j l} \psi_ {i m}) (1 \leq k, l, m \leq r).\tag{25.17}
\]

引入矩阵 \((\alpha_{ij})\)（它与 \((\psi_{ij})\) 逆步）以及无穷小变换

\[
X _ {k} = \sum_ {i = 1} ^ {n} \xi_ {k i} \frac {\partial}{\partial x _ {i}}, A _ {k} = \sum_ {j = 1} ^ {r} \alpha_ {k j} \frac {\partial}{\partial a _ {j}} (1 \leq k \leq r),\tag{25.18}
\]

(25.16) 与 (25.17) 可以分别写成：

\[
[ X _ {i}, X _ {j} ] = \sum_ {k = 1} ^ {r} c _ {i j} ^ {k} X _ {k} (1 \leq i, j \leq r)\tag{25.19}
\]

\[
\left[ A _ {i}, A _ {j} \right] = \sum_ {k = 1} ^ {r} c _ {i j} ^ {k} A _ {k}.\tag{25.20}
\]

反过来，若给定 r 个无穷小变换 \(X_{k}(1 \leq k \leq r)\)，它们线性无关且满足条件 (25.19)，则由这些变换生成的单参数子群生成一个具有 r 个本质参数的变换群。

最后，第三条定理（{[}203{]}，v. I, pp.~170 与 297 以及 v. 3, p.~597）把满足 (25.19) 的无穷小变换组 \((X_{k})_{1\leq k\leq r}\) 的确定化归为一个纯代数问题：我们必须有

\[
c _ {i j} ^ {k} + c _ {j i} ^ {k} = 0\tag{25.21}
\]

\[
\sum_ {l = 1} ^ {r} (c _ {i l} ^ {m} c _ {j k} ^ {l} + c _ {k l} ^ {m} c _ {i j} ^ {l} + c _ {j l} ^ {m} c _ {k i} ^ {l}) = 0 (1 \leq i, j, k, m \leq r).\tag{25.22}
\]

反过来，\(^{10}\) 若 (25.21) 与 (25.22) 得到满足，就存在一个满足关系式 (25.19) 的无穷小变换组，由此得到一个具有 r 个参数的变换群（换句话说，诸 \(X_{k}\) 的常系数线性组合组成一个李代数，而且反过来每个有限维李代数都能以这种方式得到）。

这些结果由对同构问题的研究加以补充。两个变换群称为相似的，如果可以通过对变量的一个可逆坐标变换和对参数的一个可逆坐标变换从一个变到另一个：从研究之初，李就在“典范参数”的定义中以自然的方式遇到过这一概念。他证明，如果借助“变量”上的一个变换，能把一个群的无穷小变换变成另一个群的无穷小变换，则两群相似（{[}203{]}，v. I, p.~329）。此事成立的一个必要条件是两群的李代数同构，李把这表述为：两群是“gleichzusammengesetzt”（同构组成）的；但这一条件并不充分，整整一章（{[}203{]}，v. I, chap.~19）专门用来求保证两群“相似”的补充条件。另一方面，置换群理论提供了两个这类群的“holoedric 同构”（全同构，即底层“抽象”群的同构）的概念；李把这一概念移植到变换群上，并证明两个这样的群“holoedric 同构”当且仅当它们的李代数同构（{[}203{]}，v. I, p.~418）。特别地，每个变换群都与它的每个参数群全同构，这表明，当人们想研究群的结构时，群所作用的“变量”无关紧要，事实上一切都化归为李代数。\(^{11}\)

同样通过与置换群理论的类比，李引入了子群、正规子群、“meriedric 同构”（满射同态）等概念，并证明它们对应于李代数的子代数、理想与满射同态；此外，他很早就遇到了“meriedric 同构”的一个特别重要的例子，即伴随表示，并且认识了它与群的中心的联系（{[}202{]}，vol.~5, Abh. III, pp.~42-75）。对于这些结果，正如对于基本定理一样，本质的工具是雅可比–克莱布什定理，它给出微分方程组的完全可积性（所谓“弗罗贝尼乌斯定理”的形式之一）；碰巧他还借助单参数群给出了它的一个新证明（{[}203{]}，v. I, chap.~6）。

传递性与本原性的概念——它们对置换群如此重要——对“有限连续”变换群也很自然地出现，李–恩格尔的专著对此作了详细研究（{[}203{]}，v. I, chap.~13 以及散见各处）；与稳定一个点的子群的关系以及齐性空间的概念也被注意到（在不采取整体观点的前提下所能注意到的那部分）（{[}203{]}，v. I, p.~425）。

最后，在 {[}203{]} 中，“辞典”通过引入导群与可解群的概念（李称之为“可积群”；这一由微分方程理论提示的术语一直沿用到 H. 外尔的工作为止）而得以完成（{[}203{]}，v. I, p.~261 以及 v. 3, pp.~678-679）；换位元与括号之间的关系其实李早在 1883 年就已注意到（{[}202{]}，vol.~5, p.~358）。

